\documentclass[11pt,a4paper]{article}
\usepackage[margin=1in]{geometry}
\usepackage[T1]{fontenc}
\usepackage[utf8]{inputenc}
\usepackage{lmodern}
\usepackage{microtype}
\usepackage{hyperref}
\usepackage{booktabs}
\usepackage{longtable}
\usepackage{array}
\usepackage{enumitem}
\usepackage{xcolor}
\usepackage{listings}

\hypersetup{
  colorlinks=true,
  linkcolor=blue!60!black,
  urlcolor=blue!60!black
}

\title{\textbf{QSH Project Structure and Module Responsibilities}\\
\large Detailed Architecture and File-Level Role Map}
\author{Internal Technical Documentation}
\date{\today}

\lstset{
  basicstyle=\ttfamily\small,
  breaklines=true,
  frame=single,
  columns=fullflexible
}

\begin{document}
\maketitle

\begin{abstract}
This document provides a file-by-file and module-by-module map of the QSH codebase. It is separate from the research report and focuses on implementation architecture, runtime paths, security/logging responsibilities, and maintenance guidance. The goal is to make the repository easier to navigate for contributors, auditors, and future refactoring work.
\end{abstract}

\tableofcontents
\newpage

\section{Repository-Level View}

\subsection{High-Level Layout}
\begin{lstlisting}
qsh/
├─ app/                     # Runtime Python code
│  ├─ main.py               # Primary CLI entrypoint
│  ├─ client.py             # Thin wrapper to main.py
│  ├─ server.py             # Thin wrapper to main.py
│  ├─ config.py             # Shared runtime configuration model
│  ├─ agent/                # Application actions (shell/file/system)
│  ├─ core/                 # Core transport/session/auth orchestration
│  ├─ crypto/               # BB84 simulation, KDF, AEAD, signatures
│  ├─ protocol/             # Handshake/messages/packet encoder
│  ├─ server/               # Legacy dispatcher/controller registry path
│  ├─ utils/                # Logging and helper utilities
│  └─ tests/                # Unit tests
├─ scripts/
│  └─ generate_qkd_figures.py
├─ report/
│  ├─ qsh_research_report.tex
│  ├─ qsh_project_structure_manual.tex
│  ├─ build_report.ps1
│  ├─ data/*.csv
│  └─ figures/*.png
├─ docs/
│  └─ protocol.md
├─ README.md
└─ USAGE_GUIDE.md
\end{lstlisting}

\subsection{Execution Reality}
There are two important layers:
\begin{itemize}[leftmargin=1.4em]
  \item \textbf{Authoritative runtime path}: \texttt{main.py} + \texttt{core/} + \texttt{protocol/handshake.py} + \texttt{crypto/} + \texttt{agent/} + \texttt{utils/logger.py}.
  \item \textbf{Legacy/scaffolding path}: \texttt{app/server/} and \texttt{protocol/router.py} still exist but are not the main path used by current CLI.
\end{itemize}

\section{Runtime Flow Map}

\subsection{Server Path}
\begin{enumerate}[leftmargin=1.4em]
  \item CLI in \texttt{app/main.py} parses arguments and builds \texttt{AppConfig}.
  \item \texttt{utils.logger.configure\_logging()} initializes console + rotating file logging.
  \item \texttt{core.server.start()} creates TCP listener.
  \item \texttt{protocol.handshake.server\_handshake()} performs HELLO + BB84/QKD + HKDF/AES setup.
  \item First encrypted action must be \texttt{auth}; enforced by \texttt{core.server}.
  \item \texttt{core.auth\_security.AuthGuard} applies throttle + lockout.
  \item Once authenticated, actions are handled via \texttt{agent.executor} and \texttt{agent.file\_transfer}.
\end{enumerate}

\subsection{Client Path}
\begin{enumerate}[leftmargin=1.4em]
  \item CLI in \texttt{app/main.py} chooses role (\texttt{shell/upload/download/session}).
  \item \texttt{core.client} opens socket and runs \texttt{client\_handshake()}.
  \item Encrypted \texttt{auth} request is sent before any privileged action.
  \item In \texttt{session} mode, one persistent connection serves many commands.
\end{enumerate}

\section{Detailed Module Responsibilities}

\subsection{\texttt{app/main.py}}
\textbf{Role}: canonical command router for all operational modes.
\begin{itemize}[leftmargin=1.4em]
  \item Defines shared CLI arguments (network, Eve, authentication, logging).
  \item Supports roles: \texttt{server}, \texttt{shell}, \texttt{session}, \texttt{upload}, \texttt{download}, \texttt{hash-password}.
  \item Configures logging early, then dispatches to server/client functions.
\end{itemize}

\subsection{\texttt{app/config.py}}
\textbf{Role}: centralized runtime configuration dataclass.
\begin{itemize}[leftmargin=1.4em]
  \item Carries network, BB84, AES/HKDF, auth hardening, Eve, and logging settings.
  \item Supports env-based defaults for password/hash/log paths.
\end{itemize}

\subsection{\texttt{app/core/}}
\textbf{Role}: orchestration layer.
\begin{itemize}[leftmargin=1.4em]
  \item \texttt{core/server.py}: authoritative server loop, auth gate, policy enforcement, audit events.
  \item \texttt{core/client.py}: one-shot operations and persistent REPL session.
  \item \texttt{core/connection.py}: framed read/write helpers for plain + secure messages.
  \item \texttt{core/auth\_security.py}: PBKDF2 hashing, verifier check, lockout/backoff state machine.
  \item \texttt{core/session.py}: session dataclass (currently lightweight, mostly structural).
\end{itemize}

\subsection{\texttt{app/crypto/}}
\textbf{Role}: cryptographic and QKD primitives.
\begin{itemize}[leftmargin=1.4em]
  \item \texttt{crypto/bb84.py}: qubit circuit construction and batched Aer simulation; includes Eve models and sifting/QBER logic.
  \item \texttt{crypto/kdf.py}: HKDF-SHA256 derivation to encryption key + session ID.
  \item \texttt{crypto/cipher.py}: AES-GCM framed encryption with sequence-AAD anti-replay check.
  \item \texttt{crypto/signatures.py}: SHA-256 helper utility (currently auxiliary).
\end{itemize}

\subsection{\texttt{app/protocol/}}
\textbf{Role}: message/handshake protocol definitions.
\begin{itemize}[leftmargin=1.4em]
  \item \texttt{protocol/handshake.py}: most critical protocol logic (HELLO, BB84 exchange, QBER acceptance, KDF transition, QKD logs).
  \item \texttt{protocol/encoder.py}: length-prefixed JSON packet codec.
  \item \texttt{protocol/messages.py}: symbolic message constants.
  \item \texttt{protocol/router.py}: legacy helper that delegates to legacy dispatcher.
\end{itemize}

\subsection{\texttt{app/agent/}}
\textbf{Role}: concrete operation handlers.
\begin{itemize}[leftmargin=1.4em]
  \item \texttt{agent/executor.py}: shell command execution via subprocess.
  \item \texttt{agent/file\_transfer.py}: base64 read/write for upload/download.
  \item \texttt{agent/system.py}: platform metadata helper (currently optional/not hot-path).
\end{itemize}

\subsection{\texttt{app/utils/}}
\textbf{Role}: support utilities.
\begin{itemize}[leftmargin=1.4em]
  \item \texttt{utils/logger.py}: structured event logging + rotating file sink.
  \item \texttt{utils/helpers.py}: path helper utility (light usage).
\end{itemize}

\subsection{\texttt{app/server/} (legacy path)}
\textbf{Role}: earlier architecture based on controller/dispatcher/registry pattern.
\begin{itemize}[leftmargin=1.4em]
  \item \texttt{server/controller.py}: alternate server loop implementation.
  \item \texttt{server/dispatcher.py}: registry-backed action routing.
  \item \texttt{server/registry.py}: handler map abstraction.
\end{itemize}
Current CLI runtime does not primarily route through this stack; these files are useful as reference or for planned refactor consolidation.

\subsection{\texttt{app/client.py} and \texttt{app/server.py}}
\textbf{Role}: thin wrapper entrypoints.
\begin{itemize}[leftmargin=1.4em]
  \item Forward execution to \texttt{main.py}.
  \item \texttt{app/server.py} auto-inserts \texttt{server} subcommand when no role is passed.
\end{itemize}

\subsection{\texttt{app/tests/}}
\textbf{Role}: focused unit validation.
\begin{itemize}[leftmargin=1.4em]
  \item \texttt{test\_bb84.py}: BB84 sifting and Eve-effect sanity checks.
  \item \texttt{test\_crypto.py}: AES-GCM transport round-trip.
  \item \texttt{test\_security.py}: hash/verify correctness and lockout behavior.
\end{itemize}

\section{Scripts and Documentation Assets}

\subsection{\texttt{scripts/generate\_qkd\_figures.py}}
\textbf{Role}: reproducible experiment runner for report plots.
\begin{itemize}[leftmargin=1.4em]
  \item Runs Monte Carlo sweeps over Eve parameters.
  \item Exports CSV datasets (\texttt{report/data}).
  \item Generates publication plots (\texttt{report/figures}).
  \item Supports \texttt{quick/full} profiles for fast iteration vs deeper runs.
\end{itemize}

\subsection{\texttt{report/}}
\begin{itemize}[leftmargin=1.4em]
  \item \texttt{qsh\_research\_report.tex}: research narrative, formulas, figures.
  \item \texttt{qsh\_project\_structure\_manual.tex}: this architecture file.
  \item \texttt{build\_report.ps1}: automated figure+PDF build pipeline.
  \item \texttt{data/*.csv}, \texttt{figures/*.png}: generated experiment artifacts.
\end{itemize}

\subsection{\texttt{README.md}, \texttt{USAGE\_GUIDE.md}, \texttt{docs/protocol.md}}
\begin{itemize}[leftmargin=1.4em]
  \item \texttt{README.md}: concise project entry, setup and quick usage.
  \item \texttt{USAGE\_GUIDE.md}: operational guide with security/logging/report workflows.
  \item \texttt{docs/protocol.md}: protocol-level notes (message flow and assumptions).
\end{itemize}

\section{File-by-File Status Table}

\renewcommand{\arraystretch}{1.2}
\begin{longtable}{>{\ttfamily}p{0.34\linewidth} p{0.14\linewidth} p{0.46\linewidth}}
\toprule
\textnormal{\textbf{File}} & \textbf{Status} & \textbf{Primary Responsibility} \\
\midrule
\endfirsthead
\toprule
\textnormal{\textbf{File}} & \textbf{Status} & \textbf{Primary Responsibility} \\
\midrule
\endhead
app/main.py & Active & Primary CLI control-plane and runtime dispatch. \\
app/config.py & Active & Central configuration model (network/crypto/auth/eve/logging). \\
app/core/server.py & Active & Main server loop, secure channel handling, auth policy enforcement. \\
app/core/client.py & Active & Client requests + persistent session REPL. \\
app/core/connection.py & Active & Packet framing and socket transport helpers. \\
app/core/auth\_security.py & Active & Password hashing, verifier validation, lockout/backoff logic. \\
app/core/session.py & Partial & Session dataclass; structural placeholder for future richer state. \\
app/protocol/handshake.py & Active & BB84 + HKDF handshake transition and QKD telemetry. \\
app/protocol/encoder.py & Active & Plain JSON packet encoder/decoder. \\
app/protocol/messages.py & Active & Message type constants. \\
app/protocol/router.py & Legacy & Old routing wrapper over legacy dispatcher path. \\
app/crypto/bb84.py & Active & Qiskit/Aer BB84 primitives, Eve simulation, sifting/QBER. \\
app/crypto/kdf.py & Active & HKDF key schedule. \\
app/crypto/cipher.py & Active & AES-GCM transport and replay/order checks. \\
app/crypto/signatures.py & Auxiliary & SHA-256 helper utility. \\
app/agent/executor.py & Active & Remote shell execution backend. \\
app/agent/file\_transfer.py & Active & File upload/download serialization backend. \\
app/agent/system.py & Auxiliary & System information helper. \\
app/utils/logger.py & Active & Structured event logging with rotation. \\
app/utils/helpers.py & Auxiliary & Parent-directory helper. \\
app/server/controller.py & Legacy & Earlier server control loop implementation. \\
app/server/dispatcher.py & Legacy & Earlier registry-based action dispatching. \\
app/server/registry.py & Legacy & Handler registry abstraction (contains duplicated definition; cleanup candidate). \\
app/client.py & Active wrapper & Convenience entrypoint proxy to \texttt{main.py}. \\
app/server.py & Active wrapper & Convenience server entrypoint proxy to \texttt{main.py}. \\
app/tests/test\_bb84.py & Active & QKD and Eve behavior tests. \\
app/tests/test\_crypto.py & Active & AEAD transport tests. \\
app/tests/test\_security.py & Active & Auth-hardening tests. \\
\bottomrule
\end{longtable}

\section{Maintenance Notes and Refactor Priorities}
\begin{enumerate}[leftmargin=1.4em]
  \item \textbf{Consolidate duplicate architecture}: either remove legacy \texttt{app/server/*} path or unify with \texttt{core/server.py}.
  \item \textbf{Resolve wrapper/package ambiguity risks}: maintain clear distinction between module wrappers and package folders.
  \item \textbf{Fix duplicated content in \texttt{app/server/registry.py}}: currently duplicated class block can be cleaned.
  \item \textbf{Promote \texttt{core/session.py}}: if richer state tracking is needed (rekey counters, QoS, per-channel stats).
  \item \textbf{Expand tests}: integration tests for lockout reset and persistent session action mix.
\end{enumerate}

\section{Conclusion}
QSH has a clear and robust active runtime centered on \texttt{main.py}, \texttt{core/}, \texttt{protocol/handshake.py}, and \texttt{crypto/}. The project also keeps legacy scaffolding modules that can either be retired or integrated in future refactors. This file provides a stable architecture reference for onboarding, auditing, and planning ``next-step'' engineering work.

\end{document}
